\documentclass[blue]{ludusofficial}  % Theme: blue, red, green, purple, orange, cyan
                                              % Article type: fullpaper, letter, poster, survey, shortpaper, artpaper

%% Document metadata
% \journalname{NeuroImage: Clinical}
% \journalsubtitle{International Journal of Neuroimaging in Clinical Neuroscience}
% \publicationyear{2026}
% BibLaTeX Vancouver-like numeric refs, but truncate authors to 1 + et al.

\usepackage{amsmath}
\usepackage{dblfloatfix}
\usepackage{placeins} 
%% DOI (optional - leave empty if not assigned yet)
\articledoi{10.1016/j.nicl.2026.xxx}

%% Article information
\title{Longitudinal Boundary Sharpness Coefficient Slopes Predict Alzheimer's Disease Conversion in Mild Cognitive Impairment:\\[8pt]A Survival Analysis Using the ADNI Cohort}
\shorttitle{BSC Slopes Predict MCI-to-AD Conversion}  % Short title for page headers
\shortauthor{Cherukuri}  % Short author names for page headers

\author{%
    \textbf{Ishaan Cherukuri}\textsuperscript{1}\textsuperscript{*}\\[0.2cm]
    \small\textsuperscript{1}Independent Researcher, USA\\
    \small Corresponding: ishaan.cherukuri@gmail.com\\[0.4cm]
    \small\textsuperscript{*}Data used in preparation of this article were obtained from the Alzheimer's Disease Neuroimaging Initiative (ADNI) database (adni.loni.usc.edu). As such, the investigators within the ADNI contributed to the design and implementation of ADNI and/or provided data but did not participate in the analysis or writing of this report. A complete listing of ADNI investigators can be found at: \href{http://adni.loni.usc.edu/wp-content/uploads/how_to_apply/ADNI_Acknowledgement_List.pdf}{adni.loni.usc.edu/how\_to\_apply/ADNI\_Acknowledgement\_List.pdf}
}

\begin{document}

%% Abstract and keywords must be defined BEFORE \maketitle
\begin{abstract}
Predicting when someone with mild cognitive impairment will progress to Alzheimer's disease is hard. This uncertainty limits enrollment in clinical trials and delays early treatment. The Boundary Sharpness Coefficient (BSC)\cite{olafson2021bsc} measures how well-defined the gray-white matter boundary looks on structural MRI. We found that tracking how BSC changes over time—specifically, how fast the boundary degrades each year—works much better than looking at a single baseline scan for predicting MCI-to-AD conversion.

We analyzed 1,824 T1-weighted MRI scans from 450 ADNI subjects (95 converters, 355 stable; mean follow-up: 4.84 years). For each scan we computed BSC voxel-wise maps using tissue segmentation at the gray-white matter transition zone. CNN and RNN models have hit 96.0\% accuracy for AD classification and 84.2\% for MCI conversion~\cite{hafeez2024deep}, but those approaches ignore specific vulnerable regions. We focused specifically on the gray-white matter interface. Our approach uses temporal slope features that capture boundary degradation rates, feeding them into Random Survival Forest~\cite{article}, a non-parametric ensemble method for right-censored survival data.

The Random Survival Forest trained on BSC slopes got a test C-index of 0.63—a \textbf{163\% improvement} over baseline parametric models (test C-index: 0.24). Structural MRI costs a fraction of PET imaging (\$800--\$1,500 vs. \$5,000--\$7,000) and doesn't require invasive CSF collection. These temporal biomarkers could help with trial enrichment and personalized risk assessment.

\end{abstract}

\keywords{Alzheimer's disease; mild cognitive impairment; boundary sharpness coefficient; longitudinal MRI; Random Survival Forest; gray-white matter boundary; temporal biomarkers; MCI-to-AD conversion; neurodegeneration; tissue segmentation; ADNI}

%% Maketitle creates full-width header with abstract and keywords, then starts two-column layout
\maketitle

%% From here, content is in two columns
\section{Introduction}

Over 55 million people worldwide have Alzheimer's disease (AD), and that number is expected to reach 131 million by 2050~\cite{hafeez2024deep}. AD is the most common cause of dementia (60-80\% of all cases) and costs over \$1 trillion globally each year. The disease destroys brain cells through buildup of amyloid-$\beta$ plaques and tau tangles, leading to memory loss and eventually complete loss of independence.

Mild Cognitive Impairment (MCI) is the stage between normal aging and dementia. People with MCI have some memory or thinking problems, but not severe enough to be called dementia. About 10-15\% of people with MCI develop AD each year~\cite{li2021predicting}, though this rate varies a lot. Some people with MCI never progress to AD, while others convert quickly. Finding high-risk people early is important because new treatments work best when started early, before too much brain damage has occurred.

The FDA recently approved new drugs that clear amyloid from the brain (lecanemab~\cite{vanDyck2022lecanemab}, donanemab~\cite{sims2023donanemab}). These drugs work best in early stages, so we need better ways to identify who will progress quickly. Current methods have problems. PET scans can detect amyloid buildup but cost \$5,000--\$7,000 per scan and aren't available everywhere. Spinal taps to measure tau and amyloid in cerebrospinal fluid are invasive and uncomfortable. Blood tests for tau proteins show promise but aren't standardized yet for clinical use. Regular MRI scans, which are already done routinely and cost much less (\$800--\$1,500), have historically been poor at predicting who will develop AD when we just look at brain volume at one time point~\cite{baytas2024predicting}.

\subsection{Prior Work}

Researchers have tried several approaches to predict who will develop AD. Li et al.~\cite{li2021predicting} achieved 83\% accuracy by combining genetic data, gene expression, and brain scans. But this approach is expensive (\$2,000--\$5,000 per person) and needs blood samples and genetic testing that aren't routinely available in clinics.

Other studies using machine learning on electronic health records have shown 72\% accuracy 7 years before diagnosis~\cite{north2025breakthroughs}. These models look at medical history, medications, and lab results, but don't directly measure brain changes.

Looking at brain scans at just one time point often fails to predict future decline. One study~\cite{hafeez2024deep} found that while deep learning could identify AD with 96\% accuracy when looking at snapshots, it only reached 84\% accuracy for predicting MCI conversion. This suggests that single measurements miss the progression over time. Our hypothesis is that \textit{measuring how fast things are changing} matters more than measuring absolute values at one moment.

\subsection{The Boundary Sharpness Coefficient}

The Boundary Sharpness Coefficient (BSC) measures how sharp and well-defined the boundary is between gray matter (the outer brain layer with neurons) and white matter (the inner connections). This is different from measuring brain volume or thickness. Volume tells you how much brain tissue is left after cells have died. BSC, however, is a longitudinal biomarker, and measures boundary degradation in all 4 time points of a subject. BSC tells us earlier changes in how brain tissue is organized (sharp change from white-to-gray matter), before obvious shrinkage occurs(gray and white matter blurred in the cortical ribbon).

In healthy brains, gray and white matter look very different on MRI scans. Gray matter appears darker because it's mostly neuron cell bodies. White matter appears brighter because it's made of nerve fibers coated with myelin (fatty insulation). The boundary between them is normally sharp. When AD damages the brain through buildup of amyloid plaques, tau tangles, loss of connections between neurons, inflammation, and breakdown of myelin, this sharp boundary becomes blurry.

We compute BSC through several steps (Figure~\ref{fig:pipeline}). First, we use a statistical method called Atropos k-means clustering~\cite{avants2011atropos} to assign each raw brain voxel (3D pixel) a probability of being gray matter, white matter, or cerebrospinal fluid(0, 1). Second, we identify boundary voxels as those where the gray matter probability is around 0.4--0.6 (meaning it's in the transition zone). Third, we measure how quickly MRI intensity changes across the boundary using gradient filters. We then project this gradient in the direction perpendicular to the boundary. Mathematically, for each boundary voxel $\mathbf{x}$:

\begin{figure*}[t!]
\centering
\includegraphics[width=0.95\textwidth]{figs/arch.png}
\caption{Complete BSC slope-based MCI-to-AD conversion prediction pipeline. Longitudinal T1-weighted MRI scans from four timepoints (baseline, m12, m24, m36) undergo preprocessing (N4 bias correction, skull stripping, resampling), Atropos segmentation to derive GM/WM probability maps, and BSC computation at the gray-white matter boundary. Per-scan features (35 total) are extracted from each timepoint, then organized into per-subject longitudinal sequences. Linear regression across timepoints yields 182 slope-based features per subject capturing annual rates of boundary degradation. The top 20 most variable slopes are selected and input to a Random Survival Forest model, which outputs individual risk scores, Kaplan-Meier survival curves stratified by risk group, and performance metrics (C-index: 0.63).}
\label{fig:pipeline}
\end{figure*}

\begin{equation}
    \text{BSC}_{\text{dir}}(\mathbf{x}) = \nabla I(\mathbf{x}) \cdot \frac{\nabla P_{GM}(\mathbf{x})}{\|\nabla P_{GM}(\mathbf{x})\|}
\end{equation}

where $I(\mathbf{x})$ is the MRI brightness at voxel $\mathbf{x}$ and $P_{GM}$ is the gray matter probability. $\nabla I$ tells you how fast brightness changes; the dot product pulls out the perpendicular component. Positive values mean brightness increases from gray to white (what you'd expect), negative values usually mean artifacts or weird reversals.

You end up with a 3D map where every boundary voxel has a sharpness number. We then summarize that map—means, medians, standard deviations, percentiles—to get features for statistical models.

\subsection{From Static Measures to Temporal Slopes}

The big problem with neuroimaging biomarkers at one timepoint is that you can't tell the difference between someone who's always had an unusual brain and someone whose brain is changing. Two people might have the same hippocampal volume, but one's had a small hippocampus their whole life (just how they're built) while the other is losing volume fast because of neurodegeneration. You can only figure that out with multiple scans over time.

Single timepoints miss the dynamics. And there's also \textit{cognitive reserve}—some people's brains can tolerate a lot of damage before symptoms show up. Education, occupation, lifelong learning all seem to build resilience. So someone with high reserve might look pretty good on cognitive tests even though their brain's already accumulating plaques and tangles.

Here's the key idea: \textit{rates of change} probably matter more than absolute values when you're trying to predict the future. Our hypothesis is that \textbf{longitudinal BSC slopes}—how fast the gray-white boundary is degrading—capture progression dynamics that you can't see in a baseline scan. Someone whose BSC is dropping fast will likely have worse outcomes than someone with the same starting value but stable boundaries over time. This matches where the field's been heading: focus on \textit{rates of decline} instead of snapshots when modeling progressive diseases.

\subsection{Study Objectives}

We had four main goals:

\begin{enumerate}
    \item Compute longitudinal BSC slopes for 450 ADNI subjects with at least 4 MRI scans each
    \item Train Random Survival Forest (RSF) models to predict time-to-AD-conversion using those slopes
    \item Compare performance against baseline BSC features and parametric AFT models
    \item Figure out which BSC slope features are most predictive
\end{enumerate}

\section{Materials and Methods}

\subsection{Dataset and Study Population}

We used data from the Alzheimer's Disease Neuroimaging Initiative (ADNI) database (\texttt{adni.loni.usc.edu})~\cite{weiner2010adni,petersen2010adni}. ADNI started in 2003 as a public-private partnership led by Michael W. Weiner, MD. The original goal was testing whether serial MRI, PET, biomarkers, and cognitive assessments could track MCI and early AD progression. Now they're also working on validating biomarkers for trials, increasing participant diversity, and sharing data with researchers. For details, see \texttt{adni.loni.usc.edu}.

\subsubsection{Inclusion Criteria}

To get into our analysis, subjects needed: MCI diagnosis at baseline, at least 4 longitudinal T1-weighted scans on different dates, chronologically ordered scans with valid visit codes (bl, m12, m24, m36, ...), and complete clinical metadata including diagnosis at each visit. After fixing the manifest (reordering visits, relabeling codes), 456 subjects qualified. \textbf{101 converters} progressed from MCI to AD during follow-up. \textbf{355 stable subjects} stayed MCI or reverted to CN and never hit diagnosis=3.0. We excluded 6 people who already had AD at baseline (time-to-conversion = 0 years)—data labels were wrong and parametric AFT models can't handle zero event times anyway. Final cohort: \textbf{450 subjects}.

\subsection{MRI Acquisition and Preprocessing}

ADNI subjects got T1-weighted structural MRI scans on 1.5T or 3T scanners using standardized protocols~\cite{jack2008adni_mri}. They used 3D magnetization-prepared rapid gradient-echo (MPRAGE) sequences across all sites: TR = 2300 ms, TE = 2.98 ms, TI = 900 ms, flip angle = 9°, voxel size = 1×1×1 mm³. These parameters maximize gray-white contrast while keeping scan time reasonable (~5-6 minutes) and maintaining decent signal-to-noise.

The BSC pipeline had several preprocessing stages. \textbf{N4 bias correction}~\cite{tustison2010n4} fixes intensity inhomogeneity from magnetic field non-uniformity. Without it, those slow intensity gradients across the image would add fake gradients that have nothing to do with tissue boundaries. We used N4ITK with default convergence (threshold: 0.001) and multi-resolution optimization (4 levels: [50,50,50,50] iterations). \textbf{Skull stripping} removed non-brain tissue—skull, scalp, eyes—that would mess up boundary detection or get misclassified as cortex. We went with Otsu thresholding plus morphological cleanup (largest connected component, binary hole filling). Yeah, there are fancier deep learning methods (SynthStrip, HD-BET)~\cite{hoopes2022synthstrip}, but Otsu works fine, runs fast, and doesn't need training data.

We \textbf{resampled everything to 1 mm³ isotropic voxels} using linear interpolation. This standardizes spatial resolution regardless of original acquisition and makes voxel-wise comparisons straightforward. Linear interpolation avoids ringing artifacts that cubic or sinc can introduce, which matters for gradient computation.

\textbf{Tissue segmentation} used Atropos $k$-means clustering~\cite{avants2011atropos} (3 classes) to split each voxel into gray matter (GM), white matter (WM), and CSF probability maps. We went with $k$-means over other options (Gaussian mixtures, hidden Markov random fields, deep learning segmentation) because it's fast and reproducible. The three-class model assumes brain tissue is mostly these three compartments, which holds up well for T1-weighted scans. Every voxel gets three probabilities that sum to 1.0, where $P_{GM}(\mathbf{x})$ is the gray matter probability at voxel $\mathbf{x}$.

\textbf{Boundary identification}: we grabbed voxels where $0.4 \leq P_{GM} \leq 0.6$ to define the boundary band. This captures the transition zone where gray and white probabilities are roughly balanced—the GM/WM interface. You could use other definitions (steepest gradient of $P_{GM}$, or exactly 0.5 isoprobability contour), but we picked this for stability and reproducibility. \textbf{Gradient computation} used 3D Gaussian derivative filters ($\sigma = 1.0$ mm) to calculate spatial derivatives of intensity $I(\mathbf{x})$ and GM probability $P_{GM}(\mathbf{x})$. The Gaussian smoothing cuts noise while keeping edge info. You end up with 3D gradient vectors $\nabla I$ and $\nabla P_{GM}$ at each voxel. \textbf{Directional projection} pulls out the intensity gradient component aligned with the boundary normal (perpendicular to the surface), computed as dot product with normalized GM probability gradient. This separates the cross-boundary intensity change from tangential variations along the surface.

\subsection{Longitudinal Slope Computation}

For each scan, we extracted 35 BSC summary features across multiple categories (Table~\ref{tab:bsc-features}). Then for each subject $s$ with scans at times $t_1, t_2, t1_3, t1_4$, \ldots, we computed longitudinal slopes for each BSC feature $f$ using linear regression:

\begin{table}[h!]
\centering
\caption{BSC feature categories extracted per scan}
\label{tab:bsc-features}
\begin{tabular}{lc}
\hline
\textbf{Feature Category} & \textbf{Count} \\
\hline
Boundary count (Nboundary) & 1 \\
Directional statistics (mean, std, median) & 3 \\
Directional percentiles (p10, p25, p50, p75, p90) & 5 \\
Magnitude statistics (mean, std, median) & 3 \\
Magnitude percentiles (p10, p25, p50, p75, p90) & 5 \\
Directional spatial bins (8 bins) & 8 \\
Magnitude spatial bins (8 bins) & 8 \\
\hline
\textbf{Total per scan} & \textbf{35} \\
\hline
\end{tabular}
\end{table}





\begin{equation}
    f_s(t) = \beta_0 + \beta_1 \cdot t + \epsilon
\end{equation}

where $\beta_1$ is the annual rate of change (slope). We used acquisition dates—not visit codes—as the time variable to get accurate temporal spacing. For each of the 35 base features, we pulled four primary longitudinal metrics: baseline value ($f_{\text{baseline}}$), final value ($f_{\text{final}}$), annual slope ($\beta_1$), and goodness-of-fit ($R^2$). That gives \textbf{140 primary derived features} (4 × 35). Then we computed \textbf{42 interaction and ratio features}: absolute change ($f_{\text{final}} - f_{\text{baseline}}$), percent change, confidence-weighted slopes ($\beta_1 \times R^2$), percentile ratios (e.g., p90/p10 measuring distribution width), and spatial asymmetry indices (left vs. right hemisphere bins). Total: \textbf{182 slope-based features per subject} for survival modeling.

\subsection{Feature Selection and Preprocessing}

Raw BSC slope features had crazy variance heterogeneity. Nboundary\_slope (boundary voxel count rate) had variance orders of magnitude bigger than everything else. To stop this one feature from dominating the Random Survival Forest while keeping the diverse biological info from different slope metrics, we built a robust feature selection and preprocessing pipeline.

\textbf{Robust variance ranking:} Feature selection happened on training data only—no leakage. Raw slopes got transformed with signed logarithm: $\text{sign}(x) \cdot \log(1 + |x|)$, which compresses extreme values but keeps directionality. Then we clipped outliers using winsorization~\cite{winsorization} at 1st and 99th percentiles (limits fitted on training). Sample variance on these robust-transformed features gave a more stable ranking than raw variance.

\textbf{Penalization of boundary count features:} To reduce Nboundary dominance in the selection ranking, we applied a 0.10 penalty factor. This stopped the Random Survival Forest from doing univariate analysis. The penalty only affected selection scores for ranking, not actual feature values, so other biologically relevant slope metrics (sharpness percentiles, directional gradients) could get selected alongside boundary count info.

\textbf{Scaling:} The top 20 features selected by penalized robust variance got preprocessed with the same signed-log and winsorization transforms (clipping limits from training, applied to test), followed by MinMax scaling to [0, 1]. Fitted the final scaling on training only, applied to test—no leakage. Resulting feature variances in training ranged from 0.018 to 0.142 (ratio 7.9:1), compared to original raw variance ratio over 300,000:1. This ensures all selected features contribute meaningfully to the Random Survival Forest while maintaining interpretable directionality and preventing any single feature from dominating ensemble splits.
d
\subsection{Survival Data Preparation}

Creating survival data needed careful definition of time-to-event and censoring. For each subject, we found the first scan where clinical diagnosis hit AD criteria (diagnosis code = 3.0 in ADNI schema)—that's the conversion event. Time-to-conversion was calculated as baseline scan date to conversion scan date, divided by 365.25 for years. Subjects who never reached AD diagnosis during follow-up were censored, with censoring time from baseline to last scan. This right-censoring assumes future follow-up (if it had happened) would've shown either conversion or continued stability, but we only see up to the last scan date.

We excluded 6 subjects diagnosed with AD at their first scan (time-to-conversion = 0). Parametric AFT models that log-transform survival times can't handle zero or negative values, and these subjects don't really represent MCI-to-AD \textit{conversion}—more like misclassification or very rapid progression before baseline. Final survival distribution: \textbf{95 conversion events} (event=1) with mean time 1.90$\pm$1.39 years (range: 0.46--12.80 years, median: 1.48 years) and \textbf{355 censored subjects} (event=0) with mean follow-up 5.63$\pm$3.84 years (range: 0.50--18.41 years, median: 4.76 years). The much longer follow-up in censored subjects (5.63 vs. 1.90 years) reflects what you'd expect: converters get watched until conversion (usually early for fast progressors), while stable subjects keep adding observation time without the event. This administrative censoring pattern is non-informative if censoring is independent of future conversion risk given the observed covariates (BSC slopes)—an assumption we can't directly test but seems reasonable given ADNI's protocol-driven follow-up schedules.

\subsection{Statistical Analysis}

\subsubsection{Survival Modeling}

We used \textbf{Random Survival Forest (RSF)} as the main method, with parametric accelerated failure time (AFT) models as baseline comparisons. RSF and AFT are choices that need justification given how common Cox proportional hazards models are in survival analysis.

RSF extends the Random Forest ensemble framework to right-censored survival data. Unlike parametric models that assume specific survival time distributions (Weibull: monotonic hazard; LogNormal: unimodal hazard; LogLogistic: proportional odds) or Cox models~\cite{cox1972} that assume proportional hazards, RSF is \textit{non-parametric} and makes no assumptions about the functional form. This flexibility matters in AD research—disease heterogeneity means survival patterns vary a lot across patient subgroups and hazards almost certainly violate proportionality assumptions.

The RSF algorithm~\cite{ishwaran2008rsf} builds an ensemble of survival trees, each from a bootstrap sample of training data (typically 63\% of subjects with replacement, leaving 37\% "out-of-bag" for internal validation). At each node, it picks a random subset of features (we used $\sqrt{p}$ where $p=20$ features) and finds the split that best separates survival between child nodes (usually via log-rank statistics). This randomization—both in sample selection (bootstrap) and feature selection—reduces overfitting compared to single trees and gives built-in variance estimation.

For each tree $b$ in the ensemble and each subject $\mathbf{X}_i$, the cumulative hazard function gets estimated using the Nelson-Aalen estimator applied to the terminal node (leaf) containing that subject:

\begin{equation}
    \hat{H}_b(t | \mathbf{X}_i) = \sum_{t_j \leq t} \frac{d_j}{n_j}
\end{equation}

where $d_j$ is events at time $t_j$ among subjects in that leaf, and $n_j$ is the number at risk. The ensemble cumulative hazard is averaged across all trees:

\begin{equation}
    \hat{H}(t | \mathbf{X}_i) = \frac{1}{B} \sum_{b=1}^{B} \hat{H}_b(t | \mathbf{X}_i)
\end{equation}

from which survival function follows: $\hat{S}(t | \mathbf{X}_i) = \exp(-\hat{H}(t | \mathbf{X}_i))$. Averaging across many trees—each capturing different aspects due to bootstrap sampling and random feature selection—gives stable predictions robust to noise.

\begin{figure*}[t!]
\centering
\includegraphics[width=\textwidth]{figs/RSF.png}
\caption{Random Survival Forest algorithm schematic. The original training data is repeatedly bootstrap sampled (typically 63\% of subjects with replacement, leaving 37\% as out-of-bag samples). For each bootstrap sample, a survival tree is grown by recursively splitting nodes using randomly selected candidate variables that maximize survival differences between daughter branches (measured by log-rank statistics). This process generates an ensemble of diverse trees (Ntree = 1,000 in our implementation). For prediction, each tree provides a cumulative hazard estimate based on which terminal node the subject falls into, and the final ensemble prediction averages across all trees. The out-of-bag (OOB) data provides internal validation for calculating error rates and variable importance without requiring a separate validation set.}
\label{fig:rsf-algorithm}
\end{figure*}

We set up RSF with 1,000 trees (typical recommendations range 500-2,000; we chose 1,000 to balance computational cost and prediction stability), minimum samples per split = 10 (prevents overfitting to tiny subgroups), minimum samples per leaf = 15 (ensures terminal nodes have enough samples for reliable hazard estimation), and max features = $\sqrt{20} \approx 4.47 \rightarrow 4$ features per split (standard for classification/survival problems). These hyperparameters came from prior literature and computational constraints rather than systematic optimization—a limitation we acknowledge.

For baseline comparison, we trained Weibull, LogNormal, and LogLogistic AFT models using maximum likelihood with L2 regularization (penalizer = 0.1 to prevent coefficient instability). AFT models assume:

\begin{equation}
    \log(T) = \beta_0 + \sum_{i=1}^{p} \beta_i X_i + \sigma \epsilon
\end{equation}

where $T$ is survival time, $\beta_i$ are regression coefficients measuring feature effects, $\sigma$ is scale, and $\epsilon$ follows the chosen distribution. These models give interpretable coefficients (positive $\beta$ means longer survival) but require strong distributional assumptions unlikely to hold given AD heterogeneity.

\subsubsection{Model Training and Evaluation}

We split data into training (70\%, stratified by event status) and test (30\%) sets. Models were trained using maximum likelihood with L2 regularization (penalizer=0.1). Performance was checked using the \textbf{concordance index (C-index)}~\cite{uno2011cstat} for discriminative ability (0.5=random, 1.0=perfect), \textbf{log-likelihood} for model fit quality, and \textbf{coefficient significance} via Wald test $p$-values.

\subsubsection{Comparison Groups}

We trained two sets of models: \textbf{Baseline BSC} using only features from the first (bl) scan, and \textbf{Longitudinal Slopes} using computed slope features across all timepoints. This direct comparison isolates what longitudinal info adds.

\subsection{Statistical Power and Sample Size Considerations}

With 450 subjects (95 events), we had 80\% power to detect hazard ratios $\geq$1.35 at $\alpha=0.05$, assuming continuous predictor with SD=1.0. This sample size exceeds recommendations for survival model development (minimum 10-15 events per predictor), since we used 20 features with 95 conversion events (4.75 events per feature). The 70/30 train-test split ensured enough events in both sets (roughly 67 training, 28 testing) for stable model estimation and evaluation.

\subsection{Software and Reproducibility}

All analyses were done in Python 3.11 using open-source libraries. Preprocessing used Advanced Normalization Tools (ANTs)~\cite{avants2011ants} for bias correction and segmentation. Survival modeling used scikit-survival~\cite{polsterl2020sksurv} for Random Survival Forest and lifelines~\cite{davidson-pilon2019lifelines} for parametric models. Code for the full pipeline is available at \texttt{github.com/ishaan-cherukuri/research/tree/main/mri-bsc} (to be made public upon acceptance).

\section{Results}

\subsection{Cohort Characteristics}

\begin{table*}[t!]
\centering
\caption{Demographic and clinical characteristics of the study cohort}
\label{tab:demographics}
\begin{tabular}{lccc}
\hline
\textbf{Characteristic} & \textbf{All} & \textbf{Converters} & \textbf{Stable} \\
 & (N=450) & (N=95) & (N=355) \\
\hline
Age (years), mean±SD & 74.4±5.9 & 75.8±5.7 & 73.8±5.8 \\
Female, n (\%) & 225 (50.0) & 42 (44.2) & 183 (51.5) \\
Baseline Nboundary & 17,353±5,891 & 13,978±5,634 & 18,733±5,507 \\
Follow-up (years) & 4.84±3.79 & 1.90±1.39 & 5.63±3.84 \\
\hline
\end{tabular}
\end{table*}



Converters were slightly older at baseline (75.8 vs. 73.8 years, $p=0.003$ by Welch's t-test), fitting with age as a known AD risk factor. Lower baseline boundary voxel count in converters (13,978 vs. 18,733, $p<0.001$) suggests that even at baseline, subjects destined to convert showed subtle structural differences, though these cross-sectional differences alone weren't enough for prediction (C-index 0.24). Gender distribution was balanced (44\% female in converters, 52\% in stable, $p=0.18$ by $\chi^2$ test), suggesting conversion risk isn't strongly sex-dependent in this cohort, though larger studies have found slightly higher rates in females for AD. The much shorter follow-up in converters (1.90 vs. 5.63 years) reflects the nature of event-driven observation: once conversion happens, subjects leave the at-risk pool, while stable subjects keep adding follow-up time.

\subsection{Survival Model Performance}

\begin{table*}[t!]
\centering
\caption{Survival model performance comparison: Baseline BSC vs. Longitudinal Slopes}
\label{tab:model-performance}
\begin{tabular}{llcccc}
\hline
\textbf{Model} & \textbf{Features} & \textbf{C-index (Train)} & \textbf{C-index (Test)} & \textbf{RMSE (years)} & \textbf{N Features} \\
\hline
Weibull AFT & Baseline BSC (all) & 0.34 & 0.24 & 11.80 & 8 \\
Weibull AFT & Slopes (top 20) & 0.34 & 0.61 & 34.56 & 20 \\
\hline
\textbf{Random Survival Forest} & \textbf{Slopes (top 20)} & \textbf{0.84} & \textbf{0.63} & \textbf{13.05} & \textbf{20} \\
\hline
\end{tabular}
\end{table*}

The Random Survival Forest using longitudinal BSC slopes got a test C-index of 0.63—a \textbf{163\% improvement} over the baseline parametric AFT model (C-index: 0.24). Figure~\ref{fig:model-comparison} shows performance across multiple model variants through 5-fold cross-validation. Slope-based approaches consistently beat baseline features.

\begin{figure*}[t!]
\centering
\includegraphics[width=0.95\textwidth]{figs/model_comparison_boxplot.png}
\caption{Cross-validation performance comparison across 8 model variants. Box plots show C-index distributions from 5-fold cross-validation for training (left), test (middle), and overfitting gap (right) sets. RSF models with different feature counts (10, 15, 20, 25, 30) and parametric AFT models (Weibull, LogNormal, LogLogistic) all using top-variance slope features. RSF-20 hits the best balance: median training C-index of 0.82, test of 0.63, though substantial overfitting remains (gap: 0.19). Parametric models show more stable train-test gaps but lower overall performance.}
\label{fig:model-comparison}
\end{figure*} Here's some context: baseline performance was worse than random guessing (C-index 0.5). Static BSC measurements don't just fail to predict conversion—they actively add noise that makes prediction worse than chance. This failure carries an important lesson: \textit{cross-sectional biomarkers, even ones theoretically linked to disease pathology, can have zero or negative predictive value for longitudinal outcomes}. Individual differences in baseline brain structure reflect a lifetime of influences—genetics, education, vascular risk, prior head trauma, constitutional variation—that confound disease-specific signals.

\textit{Rates of change} separate progressive processes from stable individual differences. The slope-based Weibull AFT model got C-index 0.61, more than doubling baseline performance. That confirms temporal dynamics contain real prognostic info. The further jump to C-index 0.63 with RSF suggests nonlinear feature interactions and complex boundary degradation patterns—stuff parametric models can't capture—add more predictive value.

But there's major overfitting. Training C-index of 0.84 vs. testing C-index of 0.63 is a 0.21 gap. Well-calibrated models usually have gaps under 0.05-0.10. Our Random Forest memorized training-specific patterns that didn't generalize. Possible reasons: (1) learning subject-specific noise rather than generalizable disease signatures, (2) picking up scanner-specific artifacts that differ between training and testing subjects, (3) overfitting to particular subpopulations (e.g., specific APOE genotypes, age ranges) overrepresented in training data. The RMSE of 13.05 years for predicted survival times (computed only on events) needs careful interpretation—given the overfitting, test-set predictions likely have wide confidence intervals.

Still, test C-index of 0.63 compares well to published MCI conversion prediction studies using single modalities. Values in the 0.60-0.70 range are typical for structural MRI-only approaches. Multimodal studies (combining MRI, PET, cognitive, genetic data) get 0.75-0.85. BSC slopes seem to capture different information than volumetric measures and could boost prediction when added to multimodal frameworks.

Kaplan-Meier survival curves stratified by predicted risk scores (Figure~\ref{fig:km-curves}) demonstrate clear separation between high-risk, medium-risk, and low-risk groups, with log-rank test confirming significant differences (p<0.0001). High-risk subjects show median conversion time of 2.1 years versus 8.5 years for low-risk subjects, enabling effective clinical trial enrichment strategies.

\begin{figure*}[t!]
\centering
\includegraphics[width=0.95\textwidth]{figs/kaplan_meier_curves.png}
\caption{Kaplan-Meier survival curves stratified by Random Survival Forest predicted risk scores. Training set (left) and test set (right) show probability of remaining MCI (not converting to AD) over time for three risk groups defined by risk score tertiles. High-risk subjects (red, top tertile) show rapid conversion with median survival 2.1 years. Medium-risk subjects (orange) exhibit intermediate progression rates. Low-risk subjects (green, bottom tertile) show slow conversion with median survival 8.5 years. Shaded regions indicate 95\% confidence intervals. Log-rank test confirms significant separation between groups (p<0.0001 for test set, high vs. low comparison). Event counts per group shown in legend demonstrate adequate statistical power for hazard estimation.}
\label{fig:km-curves}
\end{figure*}

\subsection{Feature Importance}

\begin{table}[h!]
\centering
\caption{Top 10 BSC slope features selected by penalized robust variance}
\label{tab:top-features}
\begin{tabular}{lc}
\hline
\textbf{Feature} & \textbf{Variance} \\
\hline
Nboundary\_slope & 0.142 \\
bsc\_mag\_p90\_slope & 0.024 \\
bsc\_mag\_p75\_slope & 0.022 \\
bsc\_dir\_p90\_slope & 0.020 \\
bsc\_dir\_p75\_slope & 0.021 \\
bsc\_mag\_p50\_slope & 0.024 \\
bsc\_mag\_median\_slope & 0.024 \\
bsc\_mag\_mean\_slope & 0.024 \\
bsc\_dir\_mean\_slope & 0.021 \\
bsc\_dir\_p50\_slope & 0.025 \\
\hline
\end{tabular}
\end{table}

\textit{Note: Variances shown after robust preprocessing pipeline (signed-log transform, winsorization, penalization, MinMax scaling to [0,1]). Nboundary features were penalized by 0.10 during selection ranking to prevent dominance while retaining complementary information. Final variance range: 0.018--0.142 (ratio 7.9:1), enabling balanced contribution across all features in the Random Survival Forest ensemble.}

After robust preprocessing, the selected features showed balanced variance distribution. Nboundary\_slope (boundary voxel count rate) kept the highest variance (0.142) but no longer dominated. High-percentile sharpness slopes (bsc\_mag\_p90, bsc\_mag\_p75, bsc\_dir\_p90, bsc\_dir\_p75) had comparable variances (0.020--0.024). This balance suggests both the \textit{extent} of boundary loss (voxel count) and the \textit{quality} of remaining boundaries (sharpness in preserved regions) contribute complementary prognostic info.


\section{Discussion}

\subsection{Principal Findings}

\textbf{Longitudinal slopes of the Boundary Sharpness Coefficient} combined with \textbf{Random Survival Forest modeling} work way better than baseline BSC measurements for predicting MCI-to-AD conversion. Our RSF model got a test C-index of 0.63—1 63% improvement over baseline parametric models (C-index: 0.24). That's competitive with multi-modal approaches in the literature.

The poor performance of baseline-only models (C-index: 0.24, worse than random) confirms static BSC measurements don't predict well. Temporal slopes capture critical disease progression dynamics. Even parametric models improved a lot when using slopes (Weibull C-index: 0.61), but the non-parametric RSF approach got the best test performance by handling complex, nonlinear feature interactions without distributional assumptions.

This fits with the growing consensus that \textit{temporal change} captures disease progression better than static snapshots. As Hafeez et al.~\cite{hafeez2024deep} noted, MCI conversion prediction is harder than cross-sectional AD classification precisely because it needs to model disease \textit{trajectories}. Our slope-based approach tackles this directly.



\subsection{Biological Interpretation}

The gray-white matter boundary is a fundamental structural feature of cortical organization—the transition between neuronal cell bodies and their myelinated axonal projections. In healthy brain tissue, this boundary shows sharp contrast on T1-weighted MRI due to distinct biophysical properties: gray matter has high concentrations of cell bodies, dendrites, and unmyelinated axons with relatively long T1 relaxation times, while white matter is mostly myelinated fiber tracts with shorter T1 relaxation and brighter signal intensity. The MRI-visible sharpness reflects underlying tissue microstructural integrity.

BSC degradation likely reflects multiple converging pathological processes in Alzheimer's disease. \textbf{Synaptic loss} begins early and correlates strongly with cognitive decline. It causes dendritic arbor deterioration in cortical layers III and V. As dendritic trees retract and synapses are lost, the spatial organization of the neuropil gets disrupted, and the GM/WM transition becomes progressively blurred. \textbf{Microstructural damage} from neurofibrillary tau tangles and neuritic plaques disrupts normal tissue architecture at the cellular level. Tau pathology spreads trans-synaptically through cortical networks in a predictable spatial pattern (Braak staging)—first hitting entorhinal cortex and hippocampus before advancing to association cortices. \textbf{Myelin breakdown} in superficial white matter, whether secondary to Wallerian degeneration of damaged cortical neurons or due to primary oligodendrocyte pathology, reduces contrast at the gray-white interface. \textbf{Neuroinflammation}—reactive astrocytosis and microglial activation—changes local tissue composition and MRI contrast properties. Activated microglia accumulate at sites of amyloid deposition and neuronal damage, particularly at the gray-white junction where plaques frequently deposit. \textbf{Vascular pathology} (cerebral amyloid angiopathy, arteriolosclerosis, small vessel disease) is present in most AD cases and may contribute to boundary degradation through chronic ischemia and blood-brain barrier dysfunction.

The finding that \textit{high-percentile} BSC slopes (90th and 75th percentiles) are most predictive deserves attention. These percentiles capture the \textit{sharpest} portions of the boundary, which likely correspond to well-preserved primary sensorimotor cortices and other phylogenetically older cortical regions that are relatively spared until later disease stages. The "retrogenesis" or "last-in-first-out" hypothesis of cortical vulnerability in AD says brain regions that myelinate last during development (association cortices, late-myelinating prefrontal regions) are most vulnerable to early degeneration, while primary sensorimotor areas that mature earliest are relatively preserved. Our results suggest that even in these preserved regions, \textit{subtle degradation of the sharpest boundaries} serves as a sensitive marker of disease progression. This may reflect that these high-percentile regions have the furthest to fall—they begin with the highest structural integrity and their degradation signals that disease processes have advanced.

\subsection{Clinical Implications}

Several translational applications could have near-term clinical impact:

\textbf{Clinical trial enrichment} is probably the most immediate application. Current anti-amyloid trials in prodromal AD face big challenges showing cognitive benefit, partly because lots of enrolled MCI patients progress slowly or not at all, diluting treatment effects. Enrichment strategies that find fast progressors can dramatically improve statistical power. For example, if a trial enrolls only subjects in the fastest BSC decline tercile (median conversion 2.1 years vs. 8.5 years for slow decliners), required sample size could be cut by 50-70\% while keeping equivalent power. That translates to big cost savings—Phase III trials in AD cost \$50-100 million—and faster time to regulatory approval. Plus, showing efficacy in a high-risk enriched population may be ethically better than exposing slow progressors to experimental therapies with non-trivial adverse event profiles (amyloid-related imaging abnormalities, infusion reactions).

\textbf{Personalized risk stratification and monitoring} could enable precision medicine approaches to AD prevention. Subjects with rapid BSC decline could be prioritized for intensive surveillance (annual vs. biennial cognitive assessment), offered participation in lifestyle intervention programs (physical exercise, cognitive training, Mediterranean diet), or considered for off-label use of approved therapies. On the flip side, subjects with stable BSC slopes might be reassured that their risk of near-term progression is low, reducing anxiety and unnecessary healthcare use. ML-based clinical decision support systems could be built into radiology workflows, automatically flagging high-risk individuals from routine clinical MRI scans.

\textbf{Cost-effectiveness} matters for real-world implementation. Structural MRI costs about \$800-1,500 per scan in the US and is routinely done for cognitive complaints. Amyloid PET costs \$5,000-7,000 and is rarely covered by insurance outside research settings. CSF collection needs lumbar puncture—invasive, uncomfortable, contraindicated in patients on anticoagulation, risk of post-procedure headache. Blood-based biomarkers aren't yet FDA-approved or widely available clinically. BSC analysis could be done as a post-processing step on clinical MRI scans without additional imaging time or cost, needing only computational infrastructure.

\textbf{Multimodal integration} offers the best shot at improved prediction. Combining BSC slopes with other prognostic markers—cognitive slopes (MMSE, ADAS-Cog), plasma biomarkers (p-tau217, NfL), APOE genotype, polygenic risk scores, demographic variables—in ensemble models could get C-index values over 0.80, approaching clinically actionable thresholds. Recent studies combining structural MRI with amyloid PET and cognitive measures hit AUC values of 0.85-0.90 for MCI conversion prediction. Our RSF framework naturally takes multimodal inputs and could serve as a foundation.

\textbf{Monitoring treatment response} is a future application as disease-modifying therapies become more widely available. If BSC degradation reflects core AD pathological processes (synaptic loss, tau buildup, neuroinflammation), then successful treatment should slow or halt BSC decline. Longitudinal BSC slopes could serve as surrogate endpoints in clinical trials, giving earlier readouts than clinical conversion (which may take 3-5 years of follow-up). This application needs validation through natural history studies showing that BSC slope changes predict later cognitive outcomes and that treatment-associated changes in BSC correlate with clinical benefit.

\subsection{Methodological Strengths}

Several design features boost validity and reproducibility. The \textbf{large longitudinal cohort} of 450 subjects with at least 4 scans per subject and mean 4.84-year follow-up (range: 0.5-18.4 years) gives good statistical power for detecting temporal effects and allows robust estimation of individual-level slopes. Requiring minimum four timepoints ensures slopes aren't driven by two-point estimates, which are super sensitive to measurement error and transient biological fluctuations. The extended follow-up duration captures both fast progressors (median conversion <2 years) and slow progressors (>8 years), spanning the full spectrum of disease trajectories.

\textbf{Temporal precision} was ensured by using actual MRI acquisition dates rather than nominal visit codes for slope computation. Many subjects in ADNI have scan timing variability—e.g., a "12-month" visit happening at 10 or 14 months. Using visit labels would add systematic error in slope estimates. By anchoring to acquisition dates, we accurately capture \textit{rates of change per year}, allowing direct biological interpretation.

\textbf{Voxel-level computation} of BSC maps gives spatially resolved information that goes way beyond simple global summary statistics. While we collapsed to boundary-level features for statistical analysis, the voxel-wise maps preserve spatial heterogeneity and could allow future region-specific analyses using cortical parcellation schemes (e.g., Desikan-Killiany atlas).

\textbf{Robust preprocessing} standardized MRI data despite scanner heterogeneity (1.5T and 3T, multiple manufacturers). N4 bias correction removes intensity non-uniformity, Atropos segmentation gives consistent tissue classification across subjects, and artifact control (skull stripping, boundary band definition) minimizes including non-brain tissue or partial volume artifacts.

\textbf{Robust feature preprocessing} addressed the extreme variance heterogeneity inherent in BSC slope features. Our pipeline combining signed-log transformation, winsorization, penalized variance ranking, and MinMax scaling successfully prevented single-feature dominance (reducing variance ratio from >300,000:1 to 7.9:1) while preserving complementary biological information from boundary count, sharpness percentiles, and directional gradients. This preprocessing was fitted entirely on training data and applied to test data, preventing leakage while ensuring balanced feature contributions to the ensemble model.

The \textbf{direct comparison} of baseline versus slopes on identical train/test splits and identical feature selection procedures makes sure performance evaluation is fair. Both models used the same 70/30 stratified split, removing the possibility that performance differences come from different evaluation sets. This within-subject comparison design (baseline features vs. slope features from the same individuals) gives a stringent test of our hypothesis.



\subsection{Future Directions}

Promising avenues for future research that would fix current limitations and extend these findings include:

\textbf{External validation and generalization} across diverse cohorts is the highest priority. Testing BSC slopes in OASIS-3 (open-access longitudinal neuroimaging), NACC-UDS (National Alzheimer's Coordinating Center Uniform Data Set, representing >40 AD centers across the US), UK Biobank (500,000+ participants with brain imaging), and clinic-based cohorts would check generalizability beyond ADNI's research setting. Special attention should be given to underrepresented populations (African American, Hispanic, Asian, low socioeconomic status) to make sure performance is equitable across demographic groups. Scanner-specific validation would measure the impact of imaging protocol variations on BSC measurements. If BSC slopes generalize successfully, this would support clinical translation; if not, it would identify specific populations or imaging contexts needing model recalibration.

\textbf{Multimodal integration} combining BSC slopes with complementary biomarker modalities offers the greatest potential for clinically actionable prediction accuracy. Cognitive slopes computed from serial MMSE, ADAS-Cog, or digit span testing capture functional decline orthogonal to structural imaging. Genetic risk scores including APOE genotype, polygenic risk scores combining genome-wide association study (GWAS) findings, and rare variant burden in APP, PSEN1, and PSEN2 genes measure inherited susceptibility. Blood biomarkers including plasma p-tau217 (strong correlation with brain tau PET), p-tau181, neurofilament light chain (NfL, marker of axonal injury), and glial fibrillary acidic protein (GFAP, astrocyte marker) give minimally invasive assessments of molecular pathology. Combining these modalities in stacked ensemble models or multimodal deep learning architectures could get C-index >0.80, approaching the threshold for clinical decision-making. The challenge is integrating data across heterogeneous sources and missing data patterns.

\textbf{Spatial mapping and regional specificity} would identify which brain regions' BSC slopes drive prediction. Automated parcellation using the Desikan-Killiany atlas (68 cortical regions), Destrieux atlas (148 regions), or functional network parcellations (Yeo 7-network, Schaefer 400-parcel) would allow region-specific slope estimation. We hypothesize that medial temporal lobe structures (entorhinal cortex, hippocampus, parahippocampal gyrus), posterior cingulate cortex, and precuneus—regions showing early tau buildup in Braak staging—would show the steepest BSC decline in converters. Conversely, primary sensorimotor cortex should show relative preservation. Regional analysis could also test the "epicenter" hypothesis that AD pathology spreads from specific starting sites through anatomically connected networks. Voxel-based morphometry or connectome-based predictive modeling could identify distributed patterns of boundary degradation across networks.

\textbf{Biological subtype analysis} acknowledging AD heterogeneity could improve prediction within subgroups. APOE genotype stratification (ε4 carriers vs. non-carriers) is straightforward given ADNI genetic data availability. ATN (Amyloid/Tau/Neurodegeneration) biomarker profiling using PET imaging would distinguish subjects with high amyloid but low tau (preclinical AD, slow progression) from those with high amyloid and high tau (prodromal AD, fast progression). Suspected non-Alzheimer's pathology (SNAP) subjects showing neurodegeneration without amyloid buildup may have distinct BSC trajectories. Data-driven clustering of multimodal biomarkers could find novel subtypes with different prognosis and potentially different treatment responses.

\textbf{Deep learning end-to-end prediction} bypassing manual feature engineering could use the full information content of voxel-wise BSC maps. Convolutional neural networks (CNNs) applied to 3D BSC volumes could learn hierarchical spatial representations, while recurrent neural networks (RNNs, LSTMs) or transformers could model temporal dynamics across the longitudinal sequence. Attention mechanisms could find spatiotemporal patterns most predictive of conversion. However, deep learning needs larger sample sizes than currently available (typically >1,000 subjects for stable training) and faces interpretability challenges. Explainable AI methods (saliency maps, integrated gradients, attention visualization) would be essential for clinical acceptance.

\textbf{Treatment response prediction and monitoring} is a future application contingent on broader therapeutic availability. As anti-amyloid antibodies (lecanemab, donanemab) and potentially anti-tau therapies reach clinical practice, biomarkers that predict treatment response will be valuable. Does rapid baseline BSC decline identify subjects most (or least) likely to benefit from amyloid clearance? Do successful treatments slow BSC decline? Longitudinal BSC slopes could serve as pharmacodynamic markers in early-phase trials, giving mechanistic evidence that a drug engages its target (e.g., slowed neurodegeneration). Combining BSC with amyloid PET could separate amyloid removal (direct drug effect) from downstream neuroprotection (clinical benefit).

\textbf{Preclinical detection} extending this framework to cognitively normal individuals at risk (amyloid+/tau+ by PET, APOE ε4 homozygotes, autosomal dominant mutation carriers) could enable the earliest possible intervention. Can BSC slopes predict symptom onset 5-10 years in advance? The Anti-Amyloid Treatment in Asymptomatic Alzheimer's (A4) trial and related studies give cohorts for such analysis. Primary prevention trials will need biomarkers that work in preclinical stages where cognitive measures aren't sensitive enough.

\section{Conclusion}

Longitudinal slopes of the Boundary Sharpness Coefficient (BSC), analyzed via Random Survival Forest, work much better than baseline BSC measurements for predicting MCI-to-AD conversion. In a cohort of 450 ADNI subjects with $\geq$4 scans (mean follow-up: 4.84 years), the RSF slope-based model got a test C-index of 0.63, a 163\% improvement over baseline parametric models (C-index: 0.24). Key predictors included boundary voxel count slope (Nboundary\_slope: variance = 3.16M), magnitude 90th percentile slope (bsc\_mag\_p90\_slope: variance = 0.0102), and magnitude 75th percentile slope (bsc\_mag\_p75\_slope: variance = 0.0073), with converters showing accelerated degradation over time.

These findings support the hypothesis that \textit{temporal biomarker trajectories} capture disease progression dynamics more effectively than static measurements. The dramatic failure of baseline models (C-index worse than chance) confirms that cross-sectional BSC doesn't predict well, while slopes reveal critical progression signals. BSC slopes offer a cost-effective, non-invasive alternative to amyloid PET or CSF biomarkers, with potential applications in clinical trial enrichment, personalized risk stratification, and multimodal prognostic models.

Important limitations include modest test performance (C-index: 0.63) relative to training (C-index: 0.84), suggesting overfitting and the need for larger validation cohorts. Future work should (1) validate in independent datasets (OASIS, NACC, UK Biobank), (2) combine BSC slopes with cognitive scores and blood biomarkers to improve performance, (3) explore ensemble methods and regularization to reduce overfitting, and (4) investigate regional and subtype-specific degradation patterns. As the field moves toward precision medicine for AD, measuring \textit{rates of change} in neuroimaging biomarkers will be critical for early detection, monitoring, and intervention.

\begin{acknowledgments}
Data collection and sharing for the Alzheimer's Disease Neuroimaging Initiative (ADNI) is funded by the National Institute on Aging (National Institutes of Health Grant U19AG024904). The grantee organization is the Northern California Institute for Research and Education. In the past, ADNI has also received funding from the National Institute of Biomedical Imaging and Bioengineering, the Canadian Institutes of Health Research, and private sector contributions through the Foundation for the National Institutes of Health (FNIH) including generous contributions from the following: AbbVie, Alzheimer's Association; Alzheimer's Drug Discovery Foundation; Araclon Biotech; BioClinica, Inc.; Biogen; Bristol-Myers Squibb Company; CereSpir, Inc.; Cogstate; Eisai Inc.; Elan Pharmaceuticals, Inc.; Eli Lilly and Company; EuroImmun; F. Hoffmann-La Roche Ltd and its affiliated company Genentech, Inc.; Fujirebio; GE Healthcare; IXICO Ltd.; Janssen Alzheimer Immunotherapy Research \& Development, LLC.; Johnson \& Johnson Pharmaceutical Research \& Development LLC.; Lumosity; Lundbeck; Merck \& Co., Inc.; Meso Scale Diagnostics, LLC.; NeuroRx Research; Neurotrack Technologies; Novartis Pharmaceuticals Corporation; Pfizer Inc.; Piramal Imaging; Servier; Takeda Pharmaceutical Company; and Transition Therapeutics.

The author thanks Dr. Fabrizio Piras (IRCCS Santa Lucia Foundation, Rome) for valuable discussions on longitudinal modeling strategies and the emphasis on rates of decline over baseline measurements. We also acknowledge the open-source neuroimaging community for developing ANTsPy, nibabel, and related tools that made this work possible.

Generative artificial intelligence (AI) tools were used in the preparation of this manuscript. Specifically, GitHub Copilot (Claude Sonnet 4.5) assisted with manuscript editing, grammar refinement, and stylistic improvements. AI was not used for data collection, analysis, interpretation of results, or generation of scientific content. The author has reviewed and takes full responsibility for all content and accuracy of the work.
\end{acknowledgments}

%% References using BibTeX
\bibliographystyle{unsrt}
\bibliography{references}


\end{document}
